%%%% INSERT_AFTER_LINE: 622
%%%% ANCHOR_TEXT: % ---- END NOTES ----
% CODE INPUT
We used four datasets: QM9, with the HOMO--LUMO gap as the target, and three assay datasets, logD, Caco-2 and hERG K$_i$. On each we fitted six representations and 20 model configurations. The 20 configurations are thirteen base models and seven variant models, both defined under Models below. Every pairing was exposed to seven noise conditions, defined in Section~\ref{sec:noise_injection}, at seven noise levels. Hyperparameters for every model are given in Additional file~\ref{app:noise:sec:hyperparameters_default}.
%%%% END

%%%% INSERT_AFTER_LINE: 1035
%%%% ANCHOR_TEXT: % ---- END NOTES ----
% CODE INPUT
Mean predicted uncertainty rises with label noise, but that alone does not say what the model attributes the rise to. We expected the aleatoric part to rise with injected label noise, and the epistemic part to stay roughly flat. Three kinds of model split their uncertainty per molecule: the full BNNs with a variance head, the heteroscedastic GP, and QRF \citep{kendall2017, Rasmussen2005, Meinshausen2006}. We excluded the other GPs. Their observation noise is one number for the whole fit, so it cannot follow individual molecules \citep{Rasmussen2005}. For each model we took the rise in mean epistemic uncertainty from no added noise to noise level 1.5, over held-out molecules. We divided it by the rise in mean aleatoric uncertainty over the same range. A ratio near 0 means the added noise went to the aleatoric part. A ratio of 1 or more means the epistemic part grew as much or more (\textbf{[FIGURE NEEDED: table of the ratio of epistemic rise to aleatoric rise, one row per model (the two full BNNs with a variance head, heteroscedastic GP, QRF) and noise condition (Gaussian, grouped-shifted), one column per representation, each entry the range over datasets on held-out molecules]}).

% CODE INPUT
The full BNNs with a variance head split the added noise as expected. Their ratio is at most 0.43 in every model, dataset, representation and noise condition. QRF does not. Both of its parts rise with noise, and on some representations the epistemic part rises faster (\textbf{[FIGURE NEEDED: aleatoric and epistemic parts against noise level, one panel per model with a per-molecule split, QM9 HOMO--LUMO gap, ECFP4, Gaussian noise]}). This holds although QRF is among the more noise-robust models. On QM9 under Gaussian noise, it ranks 4th of the 19 models and variants on \aucnorm on four of the six representations. Its clean $R^2$ on the same data ranks between 6th and 13th of the 19, depending on the representation.

% CODE INPUT
Representation explained little of the variance in noise robustness, but for QRF it decides how the added noise is split. Over datasets and every noise condition except censoring, the QRF ratio on held-out molecules ranges from 0.63 to 1.09 on the three fingerprints, ECFP4, Avalon and Sort \& Slice. On PDV the same range is 1.03 to 1.42. On ChemBERTa and MHG-GNN it is 1.28 to 2.05. Under Gaussian noise alone, the held-out ranges are 0.65--0.75 on ECFP4, 0.69--0.76 on Sort \& Slice and 0.81--0.89 on Avalon. They are 1.17--1.32 on PDV, 1.54--1.72 on MHG-GNN and 1.60--1.88 on ChemBERTa. On the learned embeddings, QRF therefore puts more of the added noise into its epistemic part than into its aleatoric part.

% CODE INPUT
The type of noise also moves the split, most clearly for the heteroscedastic GP. Under Gaussian noise its ratio on held-out molecules lies between $-0.05$ and 0.36 on every representation. Under grouped-shifted noise the range over datasets is 0.73--1.00 on ECFP4, 0.44--0.61 on Sort \& Slice and 0.32--0.46 on Avalon. On PDV, ChemBERTa and MHG-GNN it is 0.13--0.22, 0.14--0.23 and 0.21--0.25. Laplace, Student-t and outlier noise give splits like those under Gaussian noise. Grouped-shifted noise raises the epistemic share of both the heteroscedastic GP and QRF, for every dataset and representation. One possible reading is that a shared offset per scaffold family looks to the model like structure in the data, not label noise.

% CODE INPUT
We next asked how closely the aleatoric part follows the amount of added noise. Held-out labels carry no injected noise, so this comparison uses out-of-fold training molecules. We fitted the slope of mean aleatoric uncertainty against the standard deviation of the added noise, both in the label's units, under Gaussian noise. We report the median over folds (\textbf{[FIGURE NEEDED: aleatoric part against the standard deviation of the added noise, per model, ECFP4 and PDV, out-of-fold training molecules, dashed line where reported noise equals added noise]}). On ECFP4 and PDV the slope runs from 0.36, for QRF on PDV on Caco-2, to 1.09, for a full BNN with a variance head on PDV on hERG K$_i$. Including ChemBERTa widens the range to 0.28, for QRF on Caco-2, and 1.62, for a full BNN with a variance head on hERG K$_i$. The aleatoric part can therefore compare noise between datasets or between noise levels. It does not estimate a dataset's own measurement error.

% CODE INPUT
Censoring behaves differently from the other conditions. Clipping labels at a limit narrows their spread. NGBoost and QRF shrink their total uncertainty with that spread on every dataset, and the GPs do the same on hERG K$_i$ and Caco-2. None of these models treats a clipped label as a noisy one. They grow more confident while the labels lose information.
%%%% END
